% ch7_e §7.12–7.13 (书页 246–254 / p260–p268)
%--D-TAIL--
%JOIN%于费米面的形状及其与区边界的距离，如 §7.5 中所讨论的那样。

应当指出，如果只有电子--声子 N 过程，那么式 (7.116) 的论证本来也适用：声子的曳引恰好抵偿电子的漂移，动量便不会有任何净的耗散。电导率将会增大，而且在没有剩余电阻时会趋于无穷。在实际金属中并未观察到这种现象，原因多半在于存在强烈的电子--声子 U 过程。

我们在这里给出的论证，也可以应用到半导体中的载流子上——如 §7.6 中那样，载流子会“曳引”与之相互作用的声子。这种效应在半导体中更大，因为只有波矢 $q$ 小的声子才散射载流子，而这些声子的平均自由程很长。在低温下，人们甚至能探测到温差电功率对样品尺寸的依赖。晶格温差电功率的符号与载流子的符号相同。

\section{Hall 效应}

让我们回到 Boltzmann 方程 (7.14)，只是现在除了电场之外还有磁场 $\mathbf{H}$。假定可以取一个弛豫时间，则该方程可写为
\begin{equation*}
e\mathbf{E}\cdot\mathbf{v}_{\mathbf{k}}\left(-\frac{\partial f^{0}}{\partial\mathcal{E}}\right)=\frac{g_{\mathbf{k}}}{\tau}+\frac{e}{\hbar c}\left(\mathbf{v}_{\mathbf{k}}\wedge\mathbf{H}\right)\cdot\frac{\partial g_{\mathbf{k}}}{\partial\mathbf{k}}.\tag{7.137}
\end{equation*}

这是关于 $\mathbf{k}$ 空间中函数 $g_{\mathbf{k}}$ 的一个微分方程。假设我们处理的是自由电子，对它们有
\begin{equation*}
\hbar\mathbf{k}=m\mathbf{v},\tag{7.138}
\end{equation*}
于是我们可以用速度而不用 $\mathbf{k}$ 来标记每个态。我们尝试如下形式的解
\begin{equation*}
g_{\mathbf{k}}=\left(-\frac{\partial f^{0}}{\partial\mathcal{E}}\right)\tau\,\mathbf{v}_{\mathbf{k}}\cdot e\mathbf{A},\tag{7.139}
\end{equation*}
其中 $\mathbf{A}$ 是一个有待确定的矢量。这当然是式 (7.19) 的类似式；在没有磁场时，$\mathbf{A}$ 结果就是 $\mathbf{E}$。

把式 (7.138) 和 (7.139) 代入式 (7.137)，得到
\begin{equation*}
\mathbf{v}\cdot\mathbf{E}=\mathbf{v}\cdot\mathbf{A}+\left(\frac{e\tau}{mc}\right)\left(\mathbf{v}\wedge\mathbf{H}\right)\cdot\mathbf{A},\tag{7.140}
\end{equation*}
显然，对 $\mathbf{v}$ 的一切值，下式都使它得到满足：
\begin{equation*}
\mathbf{E}=\mathbf{A}+\frac{e\tau}{mc}\,\mathbf{H}\wedge\mathbf{A}.\tag{7.141}
\end{equation*}

这是一个矢量方程，可以从中解出 $\mathbf{A}$。不难用初等几何证明（图 133），它的一个解是
\begin{equation*}
\mathbf{A}=\frac{\mathbf{E}-\dfrac{e\tau}{mc}\,\mathbf{H}\wedge\mathbf{E}}{1+\dfrac{e^{2}\tau^{2}}{m^{2}c^{2}}\,H^{2}}.\tag{7.142}
\end{equation*}
%FIG{133}: {（原书此图无图注。）}

但我们也可以直接从式 (7.141) 出发。取由式 (7.139) 给出的 $g_{\mathbf{k}}$，便容易看出电流就是简单地
\begin{equation*}
\mathbf{J}=\sigma_{0}\mathbf{A},\tag{7.143}
\end{equation*}
其中 $\sigma_{0}$ 是金属或半导体在没有磁场时的普通电导率；这由式 (7.19) 至 (7.23) 把 $\mathbf{E}$ 换成 $\mathbf{A}$ 即可得到。于是，由式 (7.141) 和 (7.143)
\begin{align*}
\mathbf{E}&=\frac{1}{\sigma_{0}}\mathbf{J}+\frac{e\tau}{mc}\,\mathbf{H}\wedge\frac{1}{\sigma_{0}}\mathbf{J}\\
&=\rho_{0}\mathbf{J}+\frac{e\tau}{mc}\,\rho_{0}\,\mathbf{H}\wedge\mathbf{J},\tag{7.144}
\end{align*}
其中 $\rho_{0}$ 是样品的普通电阻率。

这个方程表明，产生电流 $\mathbf{J}$ 所需的电场有两个分量。沿 $\mathbf{J}$ 的方向，即沿金属丝或金属带的长度方向，
\begin{equation*}
E_{\parallel}=\rho_{\parallel}J.\tag{7.145}
\end{equation*}
这告诉我们，样品的表观电阻不因磁场而改变；也就是说，\emph{没有磁致电阻（magneto-resistance）}。

但如果 $\mathbf{H}$ 垂直于 $\mathbf{J}$ 加上，就会有一个横向电场
\begin{equation*}
E_{H}=\frac{e\tau}{mc}\,\rho_{0}HJ.\tag{7.146}
\end{equation*}

这就是 \emph{Hall 效应（Hall effect）}。对自由电子，Hall 系数（Hall coefficient）为
\begin{align*}
R&=\frac{e\tau}{mc}\,\rho_{0}\\
&=\frac{1}{nec},\tag{7.147}
\end{align*}
这里我们用式 (7.33) 消去了弛豫时间。

这个结果——Hall 常数反比于载流子密度——用分子运动论的论证很容易理解。假定像式 (7.31) 中那样，载流子以速度 $\delta\mathbf{v}$ 漂移，那么它们将受到一个横越磁场方向的 Lorentz 力
\begin{equation*}
\frac{e}{c}\,\delta\mathbf{v}\wedge\mathbf{H}\tag{7.148}
\end{equation*}
的作用。恰好能够补偿这一侧向偏转的，正是一个横向电场
\begin{align*}
\mathbf{E}_{H}&=\frac{1}{c}\,\mathbf{H}\wedge\delta\mathbf{v}\\
&=\frac{1}{c}\,\mathbf{H}\wedge\frac{1}{ne}\mathbf{J}\tag{7.149}
\end{align*}
与 $1/n$ 成正比的原因在于：电流给定时，载流子密度越小，每个载流子就必须运动得越快，从而被磁场偏转得越厉害。

这个分子运动论的论证也告诉我们，为什么在这个简单情形下没有磁致电阻。横向的两个偏转效应恰好平衡，所以载流子在纵向电场的作用下沿金属丝行进时，就好像根本不存在磁场一样。

结果 (7.147) 实际上对任何球形的载流子集合都成立。假设我们在整个 费米面上写
\begin{equation*}
\hbar k_{F}=m^{*}v_{F}\tag{7.150}
\end{equation*}
这不过是用这个关系来定义参量 $m^{*}$。于是我们应当得到
\begin{equation*}
R=\frac{e\tau}{m^{*}c\sigma_{0}}\tag{7.151}
\end{equation*}
正如式 (7.146) 中那样。但这时（参照式 (7.25)）
\begin{align*}
\sigma_{0}&=\frac{e^{2}}{12\pi^{3}\hbar}\int\tau v_{F}\,\mathrm{d}S_{F}\\
&=\frac{e^{2}\tau}{12\pi^{3}}\int\frac{k_{F}\,\mathrm{d}S_{F}}{m^{*}}\\
&=\frac{e^{2}\tau}{m^{*}}\,\frac{1}{4\pi^{3}}\,\frac{4\pi}{3}\,k_{F}^{3}\\
&=\frac{ne^{2}\tau}{m^{*}},\tag{7.152}
\end{align*}
其中 $n$ 正是半径为 $k_{F}$ 的球内所包含的载流子数目。由式 (7.151) 和 (7.152) 便得到式 (7.147)，与 $\tau$ 和 $m^{*}$ 都无关。
%FIG{134}: {（a）“电子”球：$m^{*}$ 为正。（b）“空穴”球：$m^{*}$ 为负。}

但请注意，假如我们有的一个 \emph{空穴球}（sphere of holes），我们就应当预期 $e$ 为正，也就是说，Hall 场的符号将与电子情形下得到的相反。这与上面的论证是一致的。假如我们把围绕球心的\emph{电子} 费米面拿来处理，式 (7.150) 仍然成立，但 $m^{*}$ 为负，因为此时电子速度指向球心。这个量在式 (7.151) 中仍是负的；然而在式 (7.152) 中，分母里出现的将是 $|m^{*}|$，因为正如式 (7.23) 中那样，我们实际计算的是
\begin{equation*}
\frac{e^{2}}{12\pi^{3}\hbar}\int\tau v_{F}^{2}\,\frac{\mathrm{d}S_{F}}{|v_{F}|}=\frac{e^{2}\tau}{12\pi^{3}}\int\frac{k_{F}\,\mathrm{d}S_{F}}{|m^{*}|}.\tag{7.153}
\end{equation*}
于是，$R$ 中将含有因子 $|m^{*}|/m^{*}=-1$，电荷的表观符号便反了过来。

这些公式即使在半导体的单一载流子带中也并不严格成立，因为 $\tau(\mathcal{E})$ 在被占据的各态之间可以变化。因此，修正因子依赖于散射机制 (7.50) 与 (7.73) 对能量的依赖关系，但它通常离 1 不远。把式 (7.147) 与式 (7.34) 结合起来，就给出 Hall 迁移率（Hall mobility）
\begin{equation*}
\mu_{H}=|R|\,\sigma c\tag{7.154}
\end{equation*}
即样品中多数载流子的 Hall 迁移率；其电荷的符号由 $R$ 的符号给出。

Hall 效应以及其他电流磁效应（galvanomagnetic phenomena）的理论，显然在很大程度上依赖于 Bloch 定理以及电子态的倒格子描述。在无序系统中，$\mathbf{k}$ 不再是好的量子数，通过 Boltzmann 方程进行的推导立刻变得可疑，我们就得转向某种更基本的理论，例如 Kubo 公式 (7.15)。然而事实是，目前还没有任何精确的理论可以取代简单的“电子球”公式 (7.147)——人们观察到，即使在电子平均自由程只有原子间距量级的情形下，这个公式在液态金属（§7.5）中也符合得非常好。不过，在局域态之间发生跳跃导电（hopping conduction）的极端情形下，是否还能观察到 Hall 效应，是很值得怀疑的。

\section{双带模型：磁致电阻}

现在假设我们有两类载流子，比如说，电子和空穴。对其中每一类单独而言，我们得到的都是同样的方程 (7.144)：
\begin{equation*}
\mathbf{E}=\frac{1}{\sigma_{1}}\mathbf{J}_{1}+\beta_{1}\,\mathbf{H}\wedge\frac{1}{\sigma_{1}}\mathbf{J}_{1},\tag{7.155}
\end{equation*}
其中
\begin{equation*}
\beta_{1}=\frac{e\tau_{1}}{m_{1}c}\tag{7.156}
\end{equation*}
而 $\sigma_{1}$ 是与第一类载流子相联系的电导率，这类载流子对电流的贡献为 $\mathbf{J}_{1}$。
%FIG{135}: {来自两个载流子带的 Hall 效应贡献。}
同样，对第二类载流子，
\begin{equation*}
\mathbf{E}=\frac{1}{\sigma_{2}}\mathbf{J}_{2}+\beta_{2}\,\mathbf{H}\wedge\frac{1}{\sigma_{2}}\mathbf{J}_{2}\tag{7.157}
\end{equation*}
表示它们的电流贡献与外加电场之间的关系。我们要求的是总电流
\begin{equation*}
\mathbf{J}=\mathbf{J}_{1}+\mathbf{J}_{2}.\tag{7.158}
\end{equation*}
对式 (7.155) 和 (7.157) 逐个使用形如式 (7.142) 的解，我们得到
\begin{equation*}
\mathbf{J}=\left(\frac{\sigma_{1}}{1+\beta_{1}^{2}H^{2}}+\frac{\sigma_{2}}{1+\beta_{2}^{2}H^{2}}\right)\mathbf{E}-\left(\frac{\sigma_{1}\beta_{1}}{1+\beta_{1}^{2}H^{2}}+\frac{\sigma_{2}\beta_{2}}{1+\beta_{2}^{2}H^{2}}\right)\mathbf{H}\wedge\mathbf{E},\tag{7.159}
\end{equation*}
这是一个相当复杂的公式\footnote{译注：原书式 (7.159) 第一括号中第二项的分母印作 $1+\beta_{2}^{2}H_{2}$，应为 $1+\beta_{2}^{2}H^{2}$。}，其几何表示见图 135。

要计算 Hall 系数，本应把这个公式反解出来，把 $\mathbf{E}$ 用 $\mathbf{J}$ 和 $\mathbf{H}\wedge\mathbf{J}$ 表出。但对弱的磁场，结果很容易得到：
\begin{align*}
R&=\frac{\sigma_{1}\beta_{1}+\sigma_{2}\beta_{2}}{(\sigma_{1}+\sigma_{2})^{2}}\\
&=\frac{\sigma_{1}^{2}R_{1}+\sigma_{2}^{2}R_{2}}{(\sigma_{1}+\sigma_{2})^{2}},\tag{7.160}
\end{align*}
其中 $R_{1}$ 和 $R_{2}$ 是相应的各类载流子单独存在时的 Hall 常数。这样，若 $R_{1}$ 与 $R_{2}$ 符号相反，我们看到的就是一个折中值。在半导体中，我们应当用相应载流子群的数目和迁移率来表出分电导率，即 $\sigma_{1}=n_{1}e\mu_{1}$，等等；这样一来，只有当某一类载流子明显占多数时，观测到的 Hall 迁移率 (7.154) 才会对应于某个确定的微观量。

磁致电阻则需要更细致的分析。我们要把 $\mathbf{J}$ 与 $\mathbf{E}$ 沿 $\mathbf{J}$ 方向的分量相比较。于是
\begin{align*}
\rho&=(\mathbf{J}\cdot\mathbf{E})/J^{2}\\
&=\frac{\dfrac{\sigma_{1}}{1+\beta_{1}^{2}H^{2}}+\dfrac{\sigma_{2}}{1+\beta_{2}^{2}H^{2}}}{\left(\dfrac{\sigma_{1}}{1+\beta_{1}^{2}H^{2}}+\dfrac{\sigma_{2}}{1+\beta_{2}^{2}H^{2}}\right)^{2}+\left(\dfrac{\sigma_{1}\beta_{1}H}{1+\beta_{1}^{2}H^{2}}+\dfrac{\sigma_{2}\beta_{2}H}{1+\beta_{2}^{2}H^{2}}\right)^{2}}.\tag{7.161}
\end{align*}
把它与没有磁场时的电阻
\begin{equation*}
\rho_{0}=\frac{1}{\sigma_{1}+\sigma_{2}}.\tag{7.162}
\end{equation*}
相比较，经过一点代数运算，就得到
\begin{equation*}
\frac{\Delta\rho}{\rho_{0}}\equiv\frac{\rho-\rho_{0}}{\rho_{0}}=\frac{\sigma_{1}\sigma_{2}(\beta_{1}-\beta_{2})^{2}H^{2}}{(\sigma_{1}+\sigma_{2})^{2}+H^{2}(\beta_{1}\sigma_{1}+\beta_{2}\sigma_{2})^{2}}.\tag{7.163}
\end{equation*}

这个公式本身并不十分重要；载流子能被当作分成两群、每群都具有如此简单性质的情形并不多见。尽管如此，它仍然展示了磁致电阻现象的主要特征。

首先，$\Delta\rho$ 本质上是正的，只有当 $\beta_{1}=\beta_{2}$ 时才为零。换句话说，当两群载流子在磁场中被偏转的量不同时——或者因为它们的质量不同、电荷不同，或者因为它们在被散射之前行经的距离不同——就不可能找到一个电场，使电流的这两个分量都保持沿同一方向行进。净电流作为各分量的矢量和，变小了；介质的表观电阻因而增大。

同样典型的是，$\Delta\rho$ 在弱场下正比于 $H^{2}$，但在强场下趋于\emph{饱和}（saturate）。不过后一效应与我们把载流子的能量面取为闭合曲面这一选择有关。正如我们在 §9.4 中将要看到的，有时可以找到一些特殊的晶向，使得 $\Delta\rho$ 并不饱和。

这类公式很容易推广到有许多不同类型载流子、各自分别对电流作出贡献的情形。换句话说，我们可以处理复杂的 费米面，其不同部分具有不同的 $\beta$ 值。因此，金属中磁致电阻的存在，就是 $\beta$ 在 费米面上多变——即有效质量取不同数值，或者也可能是 $\tau$ 有变化——的证据。但实际的公式是一个复杂的积分，其中含有 $\mathcal{E}(\mathbf{k})$ 在各处的各种局域导数，既不透明，也不便于用来从观测到的 $\Delta\rho$ 值推算 费米面的形状。

最简单的情形是半导体中载流子的一个“谷”（valley）（§6.6）；假如假定 $\tau$ 在每个能量面上为常数，这类载流子对电流磁效应的贡献可以相当容易地计算出来。式 (7.137) 中的 Lorentz 力算符显然是沿既垂直于 $\mathbf{v}$ 又垂直于 $\mathbf{H}$ 的方向取对 $\mathbf{k}$ 的导数。当电流沿谷轴方向时，形如式 (7.139) 的公式中的“Hall 质量”显然就是这个“横向”方向上的系数，比如说式 (6.44) 中的 $m_{i}$。但是，取不同取向的各个谷、连同逆质量张量 (6.48) 的各种非对角分量的贡献，
加起来，得到一个单一的 Hall 系数；假如晶体具有宏观立方对称性，这个系数必定是各向同性的。然而，这一对称性限制并不适用于磁致电阻：当晶体相对于电场和磁场的方向转动时，磁致电阻可以随之改变，从而给出关于半导体电子能带结构的重要信息。

上面的计算是针对\emph{横向磁致电阻}（transverse magneto-resistance）的——磁场垂直于电流方向。当磁场与载流导线平行时，还可以观察到\emph{纵向磁致电阻}（longitudinal magneto-resistance）。我们这个简单的双带模型不给出纵向磁致电阻，因为它假定每个能带都具有球对称性。纵向效应的公式要复杂得多，因为它们要求非球形的 费米面。在金属和半导体中也观察到这种效应。事实上，单晶样品电阻的改变，是电流方向和磁场方向相对晶体轴的取向的一个复杂函数。

在原理上不难、在实践上也并非不可能设想出其他依赖于环境磁场的输运“效应”。例如，热量沿温度梯度流动时，会产生一个类似于 Hall 效应的横向 \emph{Nernst 电场}（Nernst electric field）。按照晶体对称性对这些效应进行分类，并借助 Onsager 关系找出它们之间的相互联系，本身就是一项可观的理论研究；但只要弛豫时间近似成立，这些系数全都可以由 Boltzmann 方程导出。

关于式 (7.163)，还有一点值得注意。假设两群载流子具有相同的 $\tau$ 值，那么 $\Delta\rho/\rho_{0}$ 就只是 $\tau H$ 的函数。但此时 $\tau$ 本身将反比于电阻率 $\rho_{0}$，于是我们可以写出
\begin{equation*}
\frac{\Delta\rho}{\rho_{0}}=F\left(\frac{H}{\rho_{0}}\right),\tag{7.164}
\end{equation*}
其中 $F$ 是一个依赖于金属本身性质的函数。

这就是所谓的 \emph{Köhler 定则}（Köhler's rule）。它告诉我们：原则上，我们可以把不同样品的磁致电阻测量结果画在同一张基本图上；或者，利用在极低温度下 $\rho_{0}$ 很小的极纯样品，来研究超强磁场可能产生的效应。

但这条定则显然只是一个近似。两种不同类型的载流子具有相同的弛豫时间，或者甚至它们的弛豫时间之比保持不变——无论它们是被杂质还是被声子散射、在高温还是在低温——这些都绝不是必然的。这个定则所说的不过是：载流子的偏转正比于磁场与两次碰撞之间的时间的乘积。对 Köhler 定则的偏离，则表明不同类型的散射机制对不同群落的载流子有不同的影响。
